\chapter{متى نكتب ومتى نقرأ: نضيف \nt{ACET}}
\chaptermark{نضيف \nt{ACET}}
\label{sec:acet}

\section{ما الذي ينقص}

لرقاقة الذاكرة، إلى جانب أسلاك العنوان والبيانات، خط آخر: تمكين الكتابة. ما دام مطفأً تُقرأ الذاكرة فقط؛ وحين يُشغَّل يُكتب ما على خطوط البيانات. ومن دون هذا الخط لا تعمل الذاكرة: إذا كانت كل قراءة تكتب شيئًا في الوقت نفسه، فإن المقروء، ومنه المقروء خطأً، يُكتب فورًا من جديد.

في الدماغ تبدو هذه المسألة أحدّ منها في الرقاقة. في الفصل~\ref{sec:glutcomputer} كانت الذاكرة مجموعة من عصبونات \nt{GLUT} تُبقيها وصلاتها الخاصة مشتعلة (الشكل~\ref{fig:bistable}). وفي الفصل~\ref{sec:dofa} رأينا أن هذه الوصلات تتعلم بقاعدة هيب أثناء العمل مباشرة. إذن المشابك نفسها تحفظ القديم وتستقبل الجديد، ولا يوجد طور كتابة مستقل، كما لا توجد ساعة. فبمَ تميّز الشبكة اللحظة التي يجب فيها حفظ ما تراه من اللحظة التي يجب فيها تذكّر ما تعرفه؟ في الفصل~\ref{sec:neurontypes} افترضنا أن \nt{ACET} هو الذي يمسك هذا الخط. وفي هذا الفصل نتحقق من الكيفية التي يجب أن يعمل بها خط كهذا، ولماذا لا غنى عنه.

\section{ثلاثة أفعال لسلك واحد}

عصبونات الأسيتيل كولين العاملة داخل الدماغ قليلة، مجموعة في بضع مجموعات صغيرة، أما مخارجها فتتفرق في القشرة كلها. هذه قناة بثّ، مثل \nt{DOPA} و\nt{NORA}. مستوى الأسيتيل كولين في القشرة مرتفع في اليقظة المنتبهة ومنخفض في النوم البطيء~\cite{Hasselmo1999}. وكما في الدوبامين، يحدد معنى الإشارة مستقبِلُ المتلقي (القضية~\ref{prop:receptor}): للأسيتيل كولين مستقبلات سريعة تفتح قناة، ومستقبلات بطيئة تطلق تفاعلات داخل الخلية. ويهمنا هنا ثلاثة أفعال.

\emph{الأول: إضعاف الوصلات الداخلية.} بيّنت التجارب على شرائح القشرة الشمية أن الأسيتيل كولين يكبت بشدة النقل عبر الوصلات بين عصبونات المنطقة الواحدة، ويكاد لا يمس المداخل الواصلة إليها من الخارج~\cite{HasselmoBower1992}.

\emph{الثاني: تقوية المداخل.} يرفع الأسيتيل كولين استثارية العصبونات وييسّر النقل عبر بعض مسارات الدخل؛ فتصير الإشارة الآتية من أعضاء الحس أعلى صوتًا من نشاط الشبكة نفسها~\cite{Hasselmo2006}.

\emph{الثالث: تيسير تغيّر الأوزان.} عند مستوى مرتفع من الأسيتيل كولين تتغير الوصلات بسهولة أكبر~\cite{Hasselmo2006}.

بالنسبة إلى المهندس، الأفعال الثلاثة عملية واحدة: التبديل من وضع ”اقرأ“ إلى وضع ”اكتب“. تكفّ الشبكة عن الإصغاء إلى نفسها، وتبدأ بالإصغاء إلى الدخل، وتحفظ ما تسمعه.

\section{شبكة تُكمل}

لنأخذ $N$ عصبونًا من عصبونات \nt{GLUT} موصولة بعضها ببعض. ولنحفظ في وصلاتها بضعة أنماط، أي مجموعات من $pN$ عصبونًا نشطة معًا، بقاعدة هيب~\cite{Hebb1949}. إذا قُدّم إلى هذه الشبكة نصف نمط، أثارت الوصلات النصف الناقص: الشبكة \emph{تُكمل} النمط من جزئه، كما كانت المجموعة في الشكل~\ref{fig:bistable} تُبقي نفسها. هذه ذاكرة ترابطية: العنوان فيها جزء من المحتوى. و\nt{GABA}، كما في الفصل~\ref{sec:gaba}، يُبقي العدد الكلي للعصبونات النشطة قرب $pN$، لكيلا يشتعل نمطان معًا.

لنرمز لمستوى الأسيتيل كولين بـ$a$، من $0$ إلى $1$، ولنكتب أفعاله الثلاثة في النموذج مباشرة:
\begin{empheq}[box=\keybox]{equation}
\tau\frac{\dd r_i}{\dd t} = -r_i + F\Bigl(\tfrac{1+a}{2}\,x_i + (1-a)\sum_j W_{ij}\,r_j - g\Bigr),
\qquad
\frac{\dd W_{ij}}{\dd t} = \eta\,a\,(r_i-p)(r_j-p) .
\label{eq:acetnet}
\end{empheq}
هنا $r_i$ تردد العصبون $i$، و$x_i$ دخله من أعضاء الحس، و$F$ خاصية نقل ذات عتبة، و$g$ التثبيط العام من \nt{GABA}، الذي يزداد حين يتجاوز عدد العصبونات النشطة $pN$. العامل $\frac{1+a}{2}$ تقوية الدخل، و$1-a$ إضعاف الوصلات الداخلية، و$\eta a$ معدل التعلّم. والمعادلة الثانية قاعدة هيب~\eqref{eq:hebb} بصيغة مناسبة للأنماط المتناثرة~\cite{Tsodyks1988}: يزداد الوزن حين يكون العصبونان كلاهما أنشط من المعتاد، وينقص حين يكون أحدهما فقط نشطًا. والجزء السالب من الأوزان، كالعادة، تقف وراءه وسائط \nt{GABA}. لا توجد ساعة: العصبونات والأوزان تتغير كلها باتصال.

\begin{figure}[t]
\centering
\begin{tikzpicture}[>=Stealth,
  nrn/.style={circle,draw,very thick,fill=blue!6,minimum size=0.8cm,inner sep=0pt,font=\small},
  acet/.style={circle,draw,very thick,double,double distance=1.2pt,fill=orange!15,minimum size=1.35cm,inner sep=0pt,font=\scriptsize},
  bc/.style={->,thick,densely dashed,orange!85!black}]
\foreach \i/\y in {1/2.4,2/1.2,3/0}{
  \node[nrn] (N\i) at (3,\y) {$r_\i$};
  \draw[->,thick,blue!60!black] (0,\y) node[left,black] {\small $x_\i$} -- (N\i);}
\node[left,align=right,font=\small] at (-0.5,1.2) {من أعضاء\\الحس};
\draw[->,thick,gray!60!black] (N1) to[bend left=30] (N2);
\draw[->,thick,gray!60!black] (N2) to[bend left=30] (N1);
\draw[->,thick,gray!60!black] (N2) to[bend left=30] (N3);
\draw[->,thick,gray!60!black] (N3) to[bend left=30] (N2);
\draw[->,thick,gray!60!black] (N1.east) .. controls (4.6,2.0) and (4.6,0.4) .. (N3.east);
\node[right,font=\small,gray!40!black,align=left] at (4.7,1.2) {الوصلات\\الداخلية $W$};
\node[acet] (A) at (1.2,-1.9) {\nt{ACET}};
\draw[bc] (A) -- (1.6,-0.15);
\node[font=\small,orange!60!black,anchor=east] at (0.95,-0.9) {$+$ الدخل};
\draw[bc] (A) -- (4.3,0.6);
\node[font=\small,orange!60!black,anchor=west] at (4.3,0.1) {$-$ الوصلات الداخلية، $+$ معدل التعلّم};
\node[right,font=\small,align=left] at (2.2,-1.9) {مستوى مرتفع: ”اكتب“\\مستوى منخفض: ”اقرأ“};
\end{tikzpicture}
\caption{\nt{ACET} خطًا لتمكين الكتابة. العصبونات $r_i$ تتلقى الدخل $x_i$ من أعضاء الحس (الأسهم الزرقاء)، ويثير بعضها بعضًا عبر الوصلات الداخلية $W$ (الرمادية) التي تُحفظ فيها الأنماط. الأسيتيل كولين (الأسهم المتقطعة) يقوّي الدخل، ويضعف الوصلات الداخلية، وييسّر تغيّر الأوزان، كما في~\eqref{eq:acetnet}.}
\label{fig:acetnet}
\end{figure}

\section{الكتابة والقراءة يعيق أحدهما الآخر}

لنختبر النموذج~\eqref{eq:acetnet} على مهمة تتعارض فيها الكتابة والقراءة فعلًا. شبكة من $200$ عصبون تحفظ خمسة أنماط، في كل منها $20$ عصبونًا نشطًا. والمطلوب الآن حفظ نمط سادس يتطابق نصفه مع أحد الأنماط القديمة: عشرة عصبونات مشتركة بينهما. وهذا يحدث دائمًا: الغرفة الجديدة تشبه غرفة مألوفة، والوجه الجديد يشبه وجهًا قديمًا.

\emph{الكتابة عند $a$ منخفض.} لنفصل أولًا الفعلين الأولين للأسيتيل كولين عن الثالث: نُبقي معدل التعلّم كاملًا، ولا يتغير إلا تقوية الدخل وإضعاف الوصلات الداخلية. عند $a=0$ يُشعل الدخل العصبونات العشرة المشتركة، وتُشعل هذه عبر الوصلات الداخلية النصف الباقي من النمط \emph{القديم}. لا يدع \nt{GABA} النمطين يشتعلان معًا، فيطرد النمطُ القديم، المدعوم بوصلاته، النمطَ الجديد الذي لا يستند إلا إلى الدخل. الشبكة لا ترى ما أمامها، بل ما تتذكره، وقاعدة هيب تكتب المرئي بإخلاص.

والنتيجة أن النمط الجديد لم يُحفظ أصلًا: من نصفه المميز لا تتذكر الشبكة شيئًا. أما القديم فقد فسد: إذا استُدعي بنصفه المميز جرّ معه في المتوسط ستة عصبونات غريبة من عشرة. عند $a=0.1$ يُحفظ النمط الجديد، لكنه يندمج بالقديم، ويكون الفساد أشد: كل منهما، إذا استُدعي بنصفه، يجرّ معه نحو ثمانية عصبونات غريبة؛ وبدءًا من $a=0.2$ تصير الكتابة نظيفة: العصبونات الزائدة عند التذكّر أقل من واحد (بالمتوسط على عشر محاكاة).

\emph{الكتابة بمعدل التعلّم الحقيقي.} لنجعل الآن الأسيتيل كولين يحدد معدل التعلّم أيضًا، كما في~\eqref{eq:acetnet}. في عرض واحد لا يُحفظ النمط الجديد كاملًا إلا عند $a\ge0.9$؛ وعند $a=0.75$ يُحفظ منه ثمانية أعشار، وعند $a\le0.5$ لا يُحفظ أبدًا.

\emph{القراءة.} لنقدّم إلى شبكة فيها خمسة أنماط مكتوبة كتابة نظيفة نصفَ نمط، ثم نزيله. عند $a\le0.2$ تُكمل الشبكة النمط كله وتُبقيه وحدها؛ وعند $a=0.3$ تُكمله كله تقريبًا؛ وعند $a\ge0.4$ تكون الوصلات الداخلية قد ضعفت أكثر مما ينبغي، فينطفئ النشاط بعد اختفاء التلميح (الشكل~\ref{fig:acetsim}).

\begin{figure}[t]
\centering
\begin{tikzpicture}[>=Stealth,x=9cm,y=3.2cm]
\draw[->,gray!70!black] (0,0) -- (1.1,0) node[right,black] {\small $a$};
\draw[->,gray!70!black] (0,0) -- (0,1.2);
\foreach \x in {0,0.2,0.4,0.6,0.8,1.0}{\draw[gray!60!black] (\x,-0.02) -- (\x,0.02); \node[below,font=\scriptsize] at (\x,0) {$\x$};}
\foreach \y in {0,0.5,1}{\draw[gray!60!black] (-0.01,\y) -- (0.01,\y); \node[left,font=\scriptsize] at (0,\y) {$\y$};}
\node[rotate=90,font=\small] at (-0.1,0.6) {نسبة النمط المستعاد};
\draw[very thick,blue!60!black] plot[mark=*,mark size=1.6pt] coordinates {(0,1.00) (0.1,1.00) (0.2,1.00) (0.3,0.96) (0.4,0.04) (0.5,0.00) (0.75,0.00) (1.0,0.00)};
\draw[very thick,red!70!black] plot[mark=*,mark size=1.6pt] coordinates {(0.1,0.00) (0.2,0.00) (0.3,0.00) (0.5,0.00) (0.75,0.80) (0.9,1.00) (1.0,1.00)};
\node[font=\small,blue!60!black,anchor=west,align=left] at (0.03,0.5) {القراءة:\\النمط\\من نصفه};
\node[font=\small,red!70!black,anchor=east,align=right] at (1.0,0.35) {الكتابة:\\نمط جديد\\دفعة واحدة};
\end{tikzpicture}
\caption{النموذج~\eqref{eq:acetnet}: $200$ عصبون، وفي كل نمط $20$ عصبونًا نشطًا، والمتوسط على عشر محاكاة. النقاط الزرقاء القراءة: أي نسبة من النمط تُستعاد من نصفه بعد إزالة التلميح. النقاط الحمراء الكتابة: أي نسبة من نمط جديد يشبه نصفه نمطًا قديمًا تُتذكَّر بعد عرض واحد، إذا كان الأسيتيل كولين يحدد معدل التعلّم أيضًا. القراءة تنجح عند $a\le0.3$، والكتابة عند $a\ge0.9$.}
\label{fig:acetsim}
\end{figure}

\begin{keyblock}
\begin{proposition}[الكتابة والقراءة تتطلبان وضعين مختلفين]
\label{prop:writeread}
في النموذج~\eqref{eq:acetnet}، حيث الوصلات نفسها تحفظ القديم وتتعلم الجديد، ويحدد الأسيتيل كولين معدل التعلّم أيضًا، لا يوجد مستوى $a$ تنجح عنده الكتابة والقراءة بالقدر نفسه. للكتابة يجب إضعاف الوصلات الداخلية، وإلا كتبت الشبكة ما تذكّرته لا الدخل؛ وللقراءة يجب تقويتها، وإلا لم تُكمل النمط من جزئه. يلزم مبدِّل، ويمكن لسلك بثّ واحد أن يؤدي هذا الدور.
\end{proposition}
\end{keyblock}

لو لم يغيّر الأسيتيل كولين معدل التعلّم، لصلح للعملين شريط ضيق $0.2\le a\le0.3$. لكن الشبكة عندئذ ستتعلم مع كل قراءة، أي ستعيد كتابة المقروء بأخطائه. أما الفكرة نفسها، أي كبت الوصلات الداخلية أثناء التعلّم، فمأخوذة من نماذج بُنيت على قياسات في الشرائح~\cite{HasselmoAndersonBower1992}. والأرقام أعلاه تخص نموذجنا المبسّط؛ أما أين تقع حدود الأوضاع في الدماغ بالضبط فغير معروف.

\section{إلى أي حد نصدّق العينين}

المبدِّل ”اكتب/اقرأ“ حالة قصوى لسؤال أعم: إلى أي حد نصدّق الدخل، وإلى أي حد نصدّق الذاكرة؟ لدى المهندسين جواب دقيق عن هذا السؤال، وهو يوضح لماذا يتحكم سلك واحد في الوضع ومعدل التعلّم معًا.

لتتابع الشبكة مقدارًا $s(t)$ يهيم ببطء: في كل ثانية يزداد تباينه بمقدار $q$. تُخبر أعضاء الحس بـ$y(t)=s(t)+$ضجيج شدته $R$. وأفضل تقدير $\hat s$ في الزمن المتصل يعطيه مرشح كالمان-بيوسي~\cite{KalmanBucy1961}:
\begin{empheq}[box=\keybox]{equation}
\frac{\dd\hat s}{\dd t} \;=\; \frac{P}{R}\,\bigl(y-\hat s\bigr),
\qquad
\frac{\dd P}{\dd t} \;=\; q-\frac{P^2}{R} ,
\label{eq:kalman}
\end{empheq}
حيث $P$ تباين خطأ التقدير، أي مدى عدم يقين الشبكة نفسها بما تعرفه. المعادلة الأولى هي دارة \foreignlanguage{english}{RC} نفسها التي في الغشاء في نموذج العصبون~\eqref{eq:glutneuron}، لكن بثابت زمني $R/P$ متغير. حين لا تكون الشبكة واثقة يكون $P$ كبيرًا، والثابت الزمني صغيرًا، فيتبع التقدير الدخل بسرعة: هذا ”اكتب“. وحين تكون واثقة يكون $P$ صغيرًا، ولا يكاد التقدير يصغي إلى الدخل، متمسكًا بالذاكرة: هذا ”اقرأ“.

في الحالة المستقرة $P=\sqrt{qR}$، والثابت الزمني للذاكرة $\sqrt{R/q}$: إذا انزاح العالم بمقدار واحد في مئة ثانية، أي $q=0.01$، وكان ضجيج أعضاء الحس $R=1$، فعلى الذاكرة أن تتذكر نحو عشر ثوانٍ.

الأهم هنا أن عددًا واحدًا $P/R$ يحدد شيئين معًا: وزن الدخل مقارنةً بالذاكرة، وسرعة تغيّر التقدير المحفوظ. وهذا بالضبط ما يفعله الأسيتيل كولين في~\eqref{eq:acetnet}. ووفق فرضية~\cite{YuDayan2005}، يخبر الأسيتيل كولين الشبكة بمقدار شبيه بـ$P$، هو عدم اليقين \emph{المتوقع}، الذي تعرفه الشبكة مسبقًا. ولا تكون هذه القناة بطيئة دائمًا: جزء من عصبونات \nt{ACET} يستجيب للمكافأة والعقاب في نحو عشرين ميلي ثانية، وبقوة تزداد كلما كان التعزيز أقل توقعًا~\cite{Hangya2015}.

\begin{keyblock}
\begin{proposition}[مقبض واحد لوزن الدخل ومعدل التعلّم]
\label{prop:kalman}
في المرشح الأمثل~\eqref{eq:kalman} يكون الوزن الذي يُمزج به الدخل مع الذاكرة، ومعدل التعلّم، مقدارًا واحدًا هو $P/R$. لذلك من المعقول التحكم فيهما بإشارة واحدة. ومع ثبات خصائص العالم تستقر هذه الإشارة عند مستوى ثابت $\sqrt{q/R}$، ولا حاجة إلى ضبطها إلا حين تتغير خصائص العالم.
\end{proposition}
\end{keyblock}

الجملة الثانية من القضية تفسر نتيجة الفصل~\ref{sec:nora}، حيث لم يتفوق ضبط معدل التعلّم بحسب المفاجأة على أفضل معدل ثابت: في عالم خصائصُ تغيّره ثابتة يكون أفضل معدل تعلّم ثابتًا فعلًا. ويظهر المكسب حين يصير العالم فجأة مختلفًا. عندئذ يجد التقدير نفسه فجأة بعيدًا عن الدخل، و$P$ لا يعلم بذلك. والفعل الصحيح أن يُرفع $P$ بحدة، فيُنتقل بذلك إلى وضع الكتابة. ووفق الفرضية التي ناقشناها في الفصل~\ref{sec:nora}، يخبر \nt{NORA} بتغيّر كهذا \emph{غير متوقع}~\cite{YuDayan2005}. ويأتي تقسيم العمل طبيعيًا: \nt{NORA} يلاحظ أن العالم تغيّر، و\nt{ACET} يُبقي الشبكة في وضع الكتابة إلى أن يُتعلَّم الموقف الجديد.

\section{النوم وضعًا للقراءة}

إذا كان المستوى المرتفع للأسيتيل كولين وضعَ الكتابة، فماذا تفعل الشبكة عند المستوى المنخفض؟ في النوم البطيء ينخفض مستوى الأسيتيل كولين، فتتحرر الوصلات الداخلية من الكبت، ولا يكاد يوجد دخل من أعضاء الحس. الشبكة في وضع القراءة بلا تلميح تعيد تشغيل ما هو مكتوب فيها. ويُظهر تسجيل النشاط أن العصبونات التي عملت معًا في النهار تنطلق معًا من جديد في النوم البطيء~\cite{WilsonMcNaughton1994}، بل بالترتيب نفسه~\cite{Skaggs1996}. ووفق فرضية~\cite{Hasselmo1999}، تنقل إعادات التشغيل هذه ما تعلّمته الشبكة بسرعة في النهار إلى الذاكرة الطويلة الأمد (وإطالة ومضات إعادة التشغيل القصيرة اصطناعيًا تحسّن الذاكرة~\cite{FernandezRuiz2019}): في النهار كتابة من الدخل، وفي الليل إعادة كتابة من الداخل. وبالنسبة إلى المهندس يشبه هذا الكتابة في مخزن مؤقت سريع نهارًا، وتفريغه في ذاكرة دائمة بطيئة ليلًا.

\section{الخلاصة}

\begin{table}[t]
\centering
\footnotesize
\setlength{\tabcolsep}{4pt}
\renewcommand{\arraystretch}{1.15}
\begin{tabular}{>{\raggedright}p{3.2cm}>{\raggedright}p{4.2cm}>{\raggedright\arraybackslash}p{6.0cm}}
\toprule
ما يفعله \nt{ACET} & كيف & ما هو معروف \\
\midrule
يضعف الوصلات الداخلية & العامل $1-a$ في~\eqref{eq:acetnet} & مقيس في الشرائح \\
يقوّي الدخل & العامل $\frac{1+a}{2}$ & رفع الاستثارية والنقل عبر مسارات الدخل مقيس \\
ييسّر التعلّم & المعدل $\eta a$ & مقيس \\
يبدّل ”اكتب/اقرأ“ & القضية~\ref{prop:writeread} & ينتج من النماذج؛ ولا قياس مباشرًا للوضعين \\
يخبر بعدم اليقين المتوقع & $P$ في~\eqref{eq:kalman} & فرضية \\
يحرر الشبكة لإعادة التشغيل في النوم & $a$ منخفض في النوم البطيء & المستوى المنخفض في النوم وإعادات التشغيل مقيسة؛ والنقل إلى الذاكرة الطويلة الأمد فرضية \\
\bottomrule
\end{tabular}
\caption{ما يضيفه \nt{ACET} إلى شبكة من \nt{GLUT} و\nt{GABA} و\nt{DOPA} و\nt{NORA}.}
\label{tab:acet}
\end{table}

\nt{DOPA} يقول للشبكة إلى أين تغيّر الأوزان، و\nt{NORA} بأي حزم تستعمل ما تعرفه، و\nt{ACET} هل تصغي إلى العالم أم إلى نفسها (الجدول~\ref{tab:acet}). هذا آخر خطوط التحكم بالبثّ في الجدول~\ref{tab:types}: تمكين الكتابة، الذي لولاه لأفسدت الذاكرةُ، التي تحفظ الأنماط في الوصلات نفسها التي تقرؤها بها، نفسَها مع كل قراءة.

\section{تمارين}

\begin{exercise}[\lvlA]
\label{ex:09:1}
\emph{كم نتذكر.} المرشح~\eqref{eq:kalman} مع $q=0.01$ و$R=1$ يتذكر نحو $10\,\text{s}$. جد $P_\infty$ والذاكرة $\sqrt{R/q}$ إذا (أ)~تضاعف ضجيج المستشعر في السعة؛ (ب)~صار العالم يهيم أسرع بمئة مرة. (ج)~كم مرة يجب أن يزداد الضجيج لكي تتذكر الشبكة دقيقة؟
\end{exercise}

\begin{exercise}[\lvlB]
\label{ex:09:2}
\emph{وضع الكتابة بعد التغيّر.} تغيّر العالم، فقفز عدم يقين المرشح~\eqref{eq:kalman} مع $q=0.01$ و$R=1$ إلى $P_0\to\infty$. (أ)~بأي ثابت زمني يعود $P$ إلى $P_\infty$؟ (ب)~بعد كم ثانية لا يتجاوز معدل التعلّم قيمته المستقرة بأكثر من $10\,\%$؟
\end{exercise}

\begin{exercise}[\lvlB]
\label{ex:09:3}
\emph{معدل تعلّم ثابت.} تتعلم الشبكة وفق $\dd\hat s/\dd t=(y-\hat s)/T$؛ والعالم يهيم، $\dd s/\dd t=w$، $y=s+n$، حيث $w$ و$n$ ضجيجان أبيضان شدتاهما $q$ و$R$. جد أفضل $T$ وتباين الخطأ عنده. كم مرة يزداد إذا أخطأنا في $T$ بمعامل $2$ و$10$؟
\end{exercise}

\begin{exercise}[\lvlB]
\label{ex:09:4}
\emph{أين حد الكتابة.} في \texttt{models/\allowbreak acet\_memory.py} العتبة $\theta=0.35$، والوزن داخل النمط $w=0.041$ (المثالي $0.045$). النمط الجديد يتقاسم $m=10$ عصبونات مع القديم. عند أي $a$ في~\eqref{eq:acetnet} لا تشتعل بقية عصبونات النمط القديم من هذه الـ$m$ (العتبة حادة)؟
\end{exercise}

\begin{exercise}[\lvlB]
\label{ex:09:5}
\emph{نمذجة: سعة الذاكرة.} عدّل في \texttt{models/\allowbreak acet\_memory.py} الدالة \texttt{read\_sweep}: عند $a=1$ اكتب $M$ نمطًا ($M$ من $5$ إلى $100$)، ثم عند $a=0$ اقرأ العشرة الأولى من أنصافها. عند أي $M$ تظهر الأخطاء، وكيف يعتمد الحد $M_{\max}$ على $N$ و$p$؟
\end{exercise}

\begin{exercise}[\lvlB]
\label{ex:09:6}
\emph{تصميم: كتابة بلا تسرّب.} مع معدل التعلّم $\eta a$ تتعلم الشبكة أثناء القراءة أيضًا. اقترح دالة رتيبة $\eta(a)\le\eta$ تساوي صفرًا عند $a\le0.3$ ولا تقل عن $\eta a$ عند $a\ge0.9$. تحقق: بعد $20$ قراءة عند $a=0.2$ بتلميح فيه ثلاثة عصبونات غريبة، كم منها دخل النمط؟
\end{exercise}

\begin{exercise}[\lvlC]
\label{ex:09:7}
\emph{كيف يُنقل المخزن المؤقت ليلًا.} (أ)~في النوم $a=0.05$، وتدوم إعادة التشغيل $0.1\,\text{s}$ مقابل $0.2\,\text{s}$ نهارًا عند $a=1$. كم إعادة تشغيل تراكم تغيّر الأوزان الناتج عن عرض نهاري واحد، بالمعدل $\eta a$، وبالدالة $\eta(a)$ من التمرين~\ref{ex:09:6}؟ (ب)~المخزن الدائم يتعلم أبطأ من المؤقت بمئة مرة ويسمع النمط $20$ مرة في الليلة، مدة كل منها $0.1\,\text{s}$. كم ليلة يلزم؟
\end{exercise}
